% Atur variabel berikut sesuai namanya

% nama
\newcommand{\name}{PRIAGUNG RAMADHAN ALFALAKHI}
\newcommand{\authorname}{Alfalakhi, Priagung Ramadhan}
\newcommand{\nickname}{Priagung Ramadhan Alfalakhi}
\newcommand{\advisor}{Dr. Ir. Djoko Purwanto, M. Eng.}
\newcommand{\coadvisor}{Muhammad Qomaruz Zaman, S.T., M.T., Ph.D.}
\newcommand{\examinerone}{examinerone}
\newcommand{\examinertwo}{examinertwo}
\newcommand{\examinerthree}{examinerthree}
\newcommand{\headofdepartment}{Ronny Mardiyanto, S.T., M.T., Ph.D.}

% identitas
\newcommand{\nrp}{5022221161}
\newcommand{\advisornip}{196512111990021002}
\newcommand{\coadvisornip}{1992202511039}
\newcommand{\examineronenip}{12345678 901234 5 678}
\newcommand{\examinertwonip}{12345678 901234 5 678}
\newcommand{\examinerthreenip}{12345678 901234 5 678}
\newcommand{\headofdepartmentnip}{198101182003121003}

% judul
\newcommand{\tatitle}{NAVIGASI ROBOT BERODA OTONOM BERBASIS LOGIKA FUZZY DENGAN FUSI SENSOR LIDAR DAN KAMERA TERMAL}
\newcommand{\engtatitle}{\emph{FUZZY LOGIC-BASED AUTONOMOUS WHEELED ROBOT NAVIGATION USING LIDAR AND THERMAL CAMERA SENSOR FUSION}}

% tempat
\newcommand{\place}{Surabaya}

% jurusan
\newcommand{\studyprogram}{Teknik Elektro}
\newcommand{\engstudyprogram}{Electrical Engineering}

% fakultas
\newcommand{\faculty}{Teknologi Elektro dan Informatika Cerdas}
\newcommand{\engfaculty}{Intelligent Electrical and Informatics Technology}

% singkatan fakultas
\newcommand{\facultyshort}{FTEIC}
\newcommand{\engfacultyshort}{ELECTICS}

% departemen
\newcommand{\department}{Teknik Elektro}
\newcommand{\engdepartment}{Electrical Engineering}

% kode mata kuliah
\newcommand{\coursecode}{EE234899}